%% ============================================================================
%% Proceedings of the AMS (PROC) version
%% Number-theoretic emphasis; DOUBLE ANONYMOUS (no author identity)
%% ============================================================================
\documentclass[11pt]{amsart}

\usepackage{amsmath,amssymb}
\usepackage{mathrsfs}
\usepackage[hidelinks,breaklinks=true]{hyperref}
\usepackage[capitalize]{cleveref}
\usepackage{booktabs}
\usepackage{enumitem}
\usepackage{microtype}
\usepackage{graphicx}
\usepackage{float}

\newtheorem{Theorem}{Theorem}[section]
\newtheorem{Lemma}[Theorem]{Lemma}
\newtheorem{Corollary}[Theorem]{Corollary}
\theoremstyle{definition}
\newtheorem{Definition}[Theorem]{Definition}
\newtheorem{Remark}[Theorem]{Remark}

\crefname{Lemma}{Lemma}{Lemmas}
\Crefname{Lemma}{Lemma}{Lemmas}
\crefname{Theorem}{Theorem}{Theorems}
\Crefname{Theorem}{Theorem}{Theorems}
\crefname{Corollary}{Corollary}{Corollaries}
\Crefname{Corollary}{Corollary}{Corollaries}
\crefname{Definition}{Definition}{Definitions}
\Crefname{Definition}{Definition}{Definitions}
\crefname{Remark}{Remark}{Remarks}
\Crefname{Remark}{Remark}{Remarks}

\newcommand{\golden}{\varphi}
\newcommand{\Mp}{M_{p}}
\newcommand{\Ztwenty}{(\mathbb{Z}/20\mathbb{Z})^{\times}}
\newcommand{\lean}[1]{\texttt{#1}}

\sloppy
\emergencystretch=3em

\title[Simultaneous identities of the golden ratio]
      {Simultaneous identities of the golden ratio and Euler's identity,\\
       formally verified in Lean~4 over the known\\
       Mersenne prime exponents}

%% Double-anonymous submission: author identity omitted for review.

\subjclass[2020]{Primary 11A41; Secondary 11B39, 11R18, 68V20}

\keywords{Golden ratio, Mersenne primes, cyclotomic fields, modular
concentration, Gelfond--Schneider theorem, formal verification, Lean 4}

\date{\today}

\begin{document}

\begin{abstract}
We present a sustained study of a single object followed across
representations, where the content lives in the equivalences between them,
not in any representation in isolation.  We prove and formally verify in
Lean~4 that the golden ratio
$\golden=(1+\sqrt5)/2$ satisfies simultaneously twelve identities,
inclusions, and identifications across five canonical structures: the
complex rotor of Euler's identity, the real trigonometric and
real-multiplicative lines, the arithmetic of Mersenne numbers
modulo~20, and the Galois group $\Ztwenty$ of
$\mathbb{Q}(\zeta_{20})/\mathbb{Q}$. The central clause is the identity
$3\golden^{\sigma(p)}=2^{p}$, with $\sigma(p)=p\lambda-\log_\golden 3$
and $\lambda=\log_\golden 2$ transcendental by Gelfond--Schneider.
Specialised to the 52 known Mersenne prime exponents (GIMPS,
1952--2024), it holds exactly for each, and every pair of Mersenne
numbers, in their binary form $2^{p}=M_{p}+1$, stands in an exact golden
ratio. The factor~3 coincides with
$M_{2}$, with the number of admissible residue classes of
$\Mp\bmod 20$, and with the numerator of $M_{2}/|\Ztwenty|=3/8$. The
verification is conducted against Mathlib4 with $0$~\texttt{sorry}, taking as external
axioms only the Gelfond--Schneider theorem and the 44 large GIMPS
primalities.
\end{abstract}

\maketitle

\section{Introduction}

\subsection{Classification under equivalence: a millennial question}
\label{ssec:classification}

The question of which discrete objects fit canonically inside which
continuum, and how the continuum organises them by its symmetry, is
older than its modern formalisation.  Vitruvius (\emph{De
Architectura}, c.~30~BC) inscribes the human body simultaneously
within a square and a circle sharing a common centre: the body itself
plays the role of \emph{modulus}, a system of an origin, a radius, an
angle, and a magnitude, mediating between the square as discrete
arithmetic substrate (four elements, four cardinal directions, the
binary doubling $1\mapsto 2\mapsto 4$) and the circle as continuous
rotational symmetry.  Cesare Cesariano's 1521 commentary on Vitruvius
records the same module geometrically through a sequence of nested
regular polygons sharing axis (the origin) and base vertex (the
radius at angle zero), from triangle to dodecagon: each polygon
discretises the circle at the angle $2\pi/n$, all fitting inside the
same continuous rotational ambient.  The polygonal sequence is a
geometric record of a structural question: given an origin, a radius,
and a magnitude, which angular discretisations admit canonical
inscription inside the circle's rotational symmetry, and in what form.

Arguably the evolution of the same ancient structural question is
taken up in two foundational modern treatises.  Euler's
\emph{Introductio in analysin infinitorum} (1748)---characterised by
Boyer at the 1950 ICM as ``the foremost textbook of modern
times''~\cite{boyer-1951}---establishes the function as the central
object of analysis and gives a systematic analytic treatment of the
exponential and trigonometric identities, including
$e^{i\theta}=\cos\theta+i\sin\theta$---an analytic rotor of origin~$0$,
radius~$1$, angle~$\theta$, and unit magnitude---whose inverse form
$i\theta=\log(\cos\theta+i\sin\theta)$ had been known to Cotes
in~1714.  Gauss's \emph{Disquisitiones Arithmeticae} (1801, drafted
from~1796) gives the systematic arithmetic counterpart, consolidating
the work of Fermat, Euler, Lagrange, and Legendre into a unified
framework---including in particular the arithmetic of the cyclotomic
roots $\zeta_{n}=e^{2\pi i/n}$, the angles $2\pi/n$ of the regular
$n$-gon treated as algebraic integers---that, in Kline's phrasing,
``began the modern theory of numbers''~\cite{kline-1972} and set the
directions of subsequent work in the subject.

\subsection{Euler 1748: the analytic-geometric substrate}
\label{ssec:euler-substrate}

The identity
\begin{equation}\label{eq:euler}
  e^{i\theta} = \cos\theta + i\sin\theta
\end{equation}
defines the one-dimensional complex rotor: cosine and sine are the
Cartesian projections of $e^{i\theta}$ onto the real and imaginary
axes.  The
unit circle $S^{1}\subset\mathbb{C}$ is the locus of $e^{i\theta}$, and
$2\pi$ is the length of one revolution.  Within $S^{1}$, the
$n$-th roots of unity $\zeta_{n} = e^{2\pi i/n}$ are the vertices of
the regular $n$-gon: the polygonal discretisation of the circle
\emph{is} the theory of roots of unity.  Cesariano's nested polygons
are the geometric trace of what \eqref{eq:euler} encodes analytically.

The identity $e^{i\pi}+1=0$ is the evaluation of \eqref{eq:euler} at
$\theta=\pi$: an exact identity in $\mathbb{C}$ in which several
independently defined objects---the analytic exponential, the
algebraic unit~$i$, the geometric constant~$\pi$, and the arithmetic
units $1$ and~$0$---coincide simultaneously, in a way not implied by
their definitions taken individually.  The content of~\eqref{eq:euler}
is this simultaneous coincidence of independent identities at a
single object.

Subsequent mathematical work over the analytic-geometric substrate
of~\eqref{eq:euler} has produced further simultaneous identities and
related results of varying form (reciprocity laws, congruences,
equivalences, motivic realisations).  Immediately after Euler, Gauss (\emph{Disquisitiones
Arithmeticae}, 1801) established structural identities connecting the
roots of unity defined by~\eqref{eq:euler} to arithmetic data modulo
primes, opening the systematic arithmetic exploitation of the
substrate.  Over the two centuries that separate Gauss from the
present, this line has multiplied through successive
nineteenth- and twentieth-century developments; an exhaustive
account is beyond the scope of this note.  Several examples
illustrate the range.  In quantum mechanics, the Feynman path
integral expresses the propagation amplitude as a sum
$\int e^{iS[\gamma]/\hbar}\mathcal{D}\gamma$ in which the Eulerian
exponential of the action governs the analytic structure of the
theory.  In harmonic analysis on locally compact abelian groups
(Pontryagin), the characters
$\chi_\xi(x)=e^{2\pi i\langle x,\xi\rangle}$ realise the duality
between a group and its dual; the parallel theory of unitary
representations of compact Lie groups proceeds through characters
valued in the unit circle.  The Hilbert--Pólya conjecture proposes
that the nontrivial zeros of $\zeta(s)$ are the eigenvalues of a
self-adjoint operator $H$, whose associated unitary group
$\{e^{itH}\}_{t\in\mathbb{R}}$ would realise the spectrum as a
one-parameter rotation; the Selberg trace formula relates the
spectrum of the Laplacian on a hyperbolic surface to its closed
geodesics through analogous exponential sums.  In arithmetic
geometry, the Weil cohomologies place the roots of unity
$\zeta_n=e^{2\pi i/n}$ at the centre of the Frobenius action on
étale cohomology, and Grothendieck's theory of motives organises
periods $\int_\gamma\omega$---including the elementary
$\oint dz/z=2\pi i$---as fundamental invariants.  More recently,
identifications of unification across arithmetic, tropical
geometry, and motivic theory have been pursued over the
$\mathbb{F}_{1}$ substrate
(Connes--Consani, 2014--2021~\cite{connes-consani-site}).

Across these developments, arithmetic content otherwise undecidable in
the sense of Gödel~\cite{godel-1931} becomes accessible through the
analytic-geometric substrate of~\eqref{eq:euler}:
Tarski~\cite{tarski-1948} showed that elementary geometry over the
real-closed field admits a decision procedure, by contrast with
elementary arithmetic.  The field-with-one-element programme initiated
by Tits~\cite{tits-1957} and developed by Manin~\cite{manin-1995} among
others pursues a geometry over a hypothetical ``field with one
element'', under which $\mathbb{Z}$ would behave as a curve and roots of
unity supply the arithmetic with $\mathbb{F}_{1}$-structure, so as to
carry methods native to geometry over finite fields into the study
of~$\mathbb{Z}$.

This aim descends from Weil's ``Rosetta Stone''.  In a 1940 letter to his
sister Simone~\cite{weil-1940}, Weil set out an analogy among three
domains---number fields, function fields over a finite field, and Riemann
surfaces---read as three columns of a single text.  In the middle column he
was able to prove the Riemann hypothesis in its form for curves over finite
fields~\cite{weil-1948}; the ambition of the field-with-one-element
programme is to endow $\mathbb{Z}$ with enough geometry that this proof, or
an analogue of it, may be transported to the number-field column and settle
the classical Riemann hypothesis.  The Rosetta Stone is itself an instance
of the classification under equivalence of~\Cref{ssec:classification}: one
arithmetic content read across representations that are not identical but
structurally the same, the decisive work lying in the translation---the
equivalences---between them.

The present note operates from a particular angle on this substrate.
In some sense closer to the modulus of \Cref{ssec:classification} than
to those modern
extensions, it uses the golden ratio $\golden=(1+\sqrt{5})/2$
together with the rotor of Euler's identity to articulate identities
between the analytic and the arithmetic through the geometric
structure of~$\golden$, exploring the relations that $\golden$
admits between $e^{i\theta}$, $\mathbb{R}$, and $\mathbb{C}$.

In what follows the golden ratio $\golden$ enters the argument as a
structural operator between discrete polygonal geometry---specifically
the pentagon and decagon---and the continuous rotational symmetry of
the unit circle.  The relations exploited here are inspired by the
proportional compass of classical artistic practice, their use of
ratios based on the human body and the much later invention of the
\emph{sectio aurea} tool: not as an artisan's empirical
approximation, but as an exact coupling between geometric magnitudes
and algebraic roots.  Grisales Herrera~\cite{grisales} establishes
one face of this coupling between $\golden$ and the integers via the
eight reduced residue classes of $\Ztwenty$; a parallel
angular-rotation framework for primes is developed by Acevedo
Agudelo~\cite{acevedo-2021}.  The present note records, in the
analytic formulation of Euler's identity~\eqref{eq:euler}, the
broader configuration of which these facts are part.

For the golden ratio, the starting point is geometric: among the
vertices $\{\zeta_{10}^{k}\}_{k=0}^{9}$ of the regular decagon
defined by~\eqref{eq:euler} at $n=10$, the sum of two conjugate
vertices equals~$\golden$,
\[
  \golden \;=\; \zeta_{10}+\zeta_{10}^{-1}, \qquad \zeta_{10}=e^{i\pi/5}
  \qquad\text{(clause (i) of \Cref{thm:coincidence-intro})}.
\]
This inscription places $\golden$ inside the structure
of~\eqref{eq:euler} as a canonical point.  Three further inscriptions
follow from this one and constitute the content of the present note.
First, the real part of the same vertex sum gives the trigonometric
identity $\golden = 2\cos(\pi/5)$ in~$\mathbb{R}$ (clause~(iii)).
Second, the multiplicative line $\mathbb{R}_{>0}$ admits a
transcendental period $\lambda_{\log}=\log_{\golden}2$ that plays the
structural role of the $2\pi$ of~\eqref{eq:euler} and closes a turn
between the algebraic objects~$\golden$ and~$2$,
\[
  \golden^{\lambda_{\log}} \;=\; 2 \qquad\text{(clause (iv))},
\]
a real-multiplicative companion to~\eqref{eq:euler}.  Third, the
cyclotomic field $\mathbb{Q}(\zeta_{20})$ where $\golden$ resides has
Galois group $\Ztwenty$ of order~$8$ (clause~(ix)), and the Mersenne
primes $\Mp$ distribute over the structural fraction
\[
  M_{2}/|\Ztwenty| \;=\; 3/8 \qquad\text{(clause (xi))}
\]
of those eight Galois classes.  These five inscriptions, together
with seven further identities, inclusions, and identifications in the
same canonical structures (clauses (ii), (v), (vi), (vii), (viii),
(x), (xii) of \Cref{thm:coincidence-intro}), are collected in the
single formal statement~\Cref{thm:coincidence-intro} below and
verified in Lean~4.  The relevance of $\golden$ at this point in the
argument is that the inscription $\golden=\zeta_{10}+\zeta_{10}^{-1}$
in~\eqref{eq:euler} propagates---trigonometrically into~$\mathbb{R}$,
multiplicatively into~$\mathbb{R}_{>0}$, and arithmetically into the
Galois group of the cyclotomic extension---producing twelve coincident
identities under a single algebraic object.

\subsection{A simultaneous identities theorem for
\texorpdfstring{$\golden$}{phi}}\label{ssec:coincidence-intro}

By a \emph{simultaneous identities theorem} we mean a statement
asserting that a mathematical object satisfies simultaneously a
finite collection of identities provable in distinct canonical
structures, with the simultaneous validity of the collection as the
content of the statement.  The corresponding statement
for~$\golden$ is the following \Cref{thm:coincidence-intro}.  Its
twelve clauses are proved in \Cref{sec:universal}
and~\Cref{sec:triple} (clauses (i)--(iii),~(ix)
in~\Cref{rem:pentagon}, clause~(iv) in~\Cref{lem:binary}, clause~(v)
in~\Cref{thm:universal}, clause~(vi)
in~\Cref{thm:concentration-main}, clause~(vii)
in~\Cref{thm:three-main}); the four remaining clauses (viii), (x),
(xi), (xii) have short proofs but each registers a structural
knotting articulated after the statement.

The individual proofs of the clauses are short---several require only
definitional manipulation or elementary modular computation.  The
substance of the statement does not reside in the difficulty of the
separate proofs but in the simultaneity of their validity under a
single~$\golden$, and in the absence of a prior reference recording
them jointly.  The situation is analogous to that of Euler's identity
$e^{i\pi}+1=0$, whose modern proof is the evaluation of a series but
whose enunciation revealed a knotting of analysis, algebra, geometry,
and arithmetic not visible in the prior literature.  The formal
verification in Lean~4 secures the simultaneity as a single theorem,
eliminating the possibility that the clauses hold separately but not
jointly.  Throughout, $\iota:\mathbb{R}\hookrightarrow\mathbb{C}$
denotes the canonical embedding and $\golden_\mathbb{C}:=\iota(\golden)$
its image in~$\mathbb{C}$ under this embedding.

\begin{Theorem}[Simultaneous identities for $\golden$ (formal version:
\Cref{thm:euler-connector})]\label{thm:coincidence-intro}
A single $\golden$ satisfies simultaneously the following twelve
identities, inclusions, and identifications, organised by the
canonical structure they inhabit:

\smallskip
\noindent\emph{In the complex rotor of Euler's identity
$\{e^{i\theta}\}\subset\mathbb{C}$:}
\begin{enumerate}[topsep=0.2em,itemsep=0.1em,leftmargin=2.3em]
  \item[\textup{(i)}] \emph{Cyclotomic inscription:}
    $\golden_\mathbb{C} = \zeta_{10}+\zeta_{10}^{-1}$, with
    $\zeta_{10}:=e^{i\pi/5}$.
  \item[\textup{(ii)}] \emph{Tenth-root closure:} $\zeta_{10}^{10}=1$.
  \item[\textup{(xii)}] \emph{Rotor discretisation:} for every
    $k\in\mathbb{N}$, $(\zeta_{10}^{k})^{10}=1$.
\end{enumerate}

\smallskip
\noindent\emph{In the real trigonometric line, with cross-identification
to the complex inscription:}
\begin{enumerate}[topsep=0.2em,itemsep=0.1em,leftmargin=2.3em]
  \item[\textup{(iii)}] \emph{Pentagonal trigonometric identity:}
    $\golden = 2\cos(\pi/5)$.
  \item[\textup{(viii)}] \emph{Real-complex identification:}
    $\iota(\golden) = \golden_\mathbb{C}$.
\end{enumerate}

\smallskip
\noindent\emph{In the real-multiplicative line $\mathbb{R}_{>0}$:}
\begin{enumerate}[topsep=0.2em,itemsep=0.1em,leftmargin=2.3em]
  \item[\textup{(iv)}] \emph{Logarithmic binary bridge:}
    $\golden^{\lambda_{\log}}=2$, where
    $\lambda_{\log}:=\log_{\golden}2$.
  \item[\textup{(v)}] \emph{Universal real-multiplicative identity:}
    $\forall p\in\mathbb{R},\;3\,\golden^{\sigma(p)}=2^{p}$, where
    $\sigma(p):=p\,\lambda_{\log}-\log_{\golden}3$.
\end{enumerate}

\smallskip
\noindent\emph{In the arithmetic of Mersenne modulo 20:}
\begin{enumerate}[topsep=0.2em,itemsep=0.1em,leftmargin=2.3em]
  \item[\textup{(vi)}] \emph{Modular concentration:} for every prime
    $p$, $\Mp\bmod 20\in\{3,7,11\}$.
  \item[\textup{(vii)}] \emph{Triple ternary identification:}
    $3 = M_{2} = |\{3,7,11\}|$.
\end{enumerate}

\smallskip
\noindent\emph{In the Galois group of the cyclotomic extension
$\mathbb{Q}(\zeta_{20})/\mathbb{Q}$, with quantitative ratio to the
arithmetic locus:}
\begin{enumerate}[topsep=0.2em,itemsep=0.1em,leftmargin=2.3em]
  \item[\textup{(ix)}] \emph{Galois order:}
    $|\Ztwenty|=8$.
  \item[\textup{(x)}] \emph{Galois sub-object inclusion:}
    $\{3,7,11\}\subseteq\Ztwenty$.
  \item[\textup{(xi)}] \emph{Quantitative ratio:}
    $|\{3,7,11\}|/|\Ztwenty|=M_{2}/8=3/8$, equivalently
    $|\{3,7,11\}|\cdot 8 = M_{2}\cdot|\Ztwenty|$.
\end{enumerate}
\end{Theorem}

No proper sub-collection of the twelve clauses constitutes the content;
the content is the simultaneous holding of all twelve under a single
$\golden$ and in relation to Euler's identity~\eqref{eq:euler}, where
the constants $\zeta_{10}=e^{i\pi/5}$, $\lambda_{\log}=\log_{\golden}2$,
and $\sigma(p)=p\,\lambda_{\log}-\log_{\golden}3$ are formally introduced
in~\Cref{sec:universal} (\Cref{def:sigma}) and~\Cref{rem:pentagon}.
The five groupings above record the five canonical structures in
which $\golden$ is independently inscribed: the complex rotor with
its tenth-root discretisation, the real trigonometric line, the
real-multiplicative line, the arithmetic of Mersenne modulo~20, and
the Galois group of the cyclotomic decagonal extension.  The
cross-identifications (clauses (viii) and (xi)) record that the
inscriptions are inscriptions of \emph{the same}~$\golden$ across
distinct substrates and that the arithmetic locus sits as a
quantitatively specific sub-object of the Galois group.  \Cref{ssec:characterization} makes this precise as a
consistency theorem: the factor~$3$, the base~$\golden$, and the slope
$\lambda_{\log}$ are each pinned by a theorem external to the central
identity, and the identity is the unique point at which they agree.

\begin{figure}[H]
  \centering
  \includegraphics[width=\textwidth]{cd_correspondance.pdf}
  \caption{The twelve simultaneous identities of~\Cref{thm:coincidence-intro}
  (clauses (i)--(xii)), together with the reciprocal turn
  $\golden^{-\lambda_{\log}}=\tfrac12$ of~\Cref{sec:spectral}, as a single
  commutative diagram.  Every directed path commutes to the universal
  identity $3\,\golden^{\sigma(p)}=2^{p}$ (green) and to the reciprocal turn
  $\tfrac12=\golden^{-\lambda_{\log}}$ (red, left); the factor $3=M_{2}$ and
  the half each appear as a convergence of several identifications under the
  single algebraic object~$\golden$.}
  \label{fig:coincidence-diagram}
\end{figure}

Clauses (vi), (vii), (x), (xi), taken in isolation, fall in the
genealogy of arithmetic identities over
$\mathbb{Q}(\sqrt{5})$ companion to Lucas's 1878 Fibonacci
congruence and to Grisales Herrera's classification of the eight
residue classes of $\Ztwenty$ into split and inert types
in~$\mathbb{Q}(\sqrt{5})$~\cite{grisales}.  They enter
\Cref{thm:coincidence-intro} as four of the twelve simultaneous clauses
of the simultaneous identities theorem, not as the global content of the
theorem.  The splitting type of each prime $\Mp$ in $\mathbb{Q}(\sqrt5)$ is
developed in~\Cref{sec:fibonacci} via Lucas's congruence, where it is
shown to be governed by the exponent $p\bmod4$ and to single out the
inert primes as those satisfying the Lucas--Lehmer divisibility
$\Mp\mid F_{\Mp+1}$.

Beyond these twelve simultaneous clauses, the same multiplicative period
$\lambda_{\log}$, read in reverse, closes a reciprocal turn
$\golden^{-\lambda_{\log}}=\tfrac12$.  This value is not a thirteenth
clause but a point already inscribed in the identity---the reciprocal of
clause~(iv), the $p=-1$ node of the universal identity~(v), and the
multiplicative-line image of the rotor's conjugate root (clause~(i))---%
developed in~\Cref{sec:spectral}, where it is also identified as the
value at which the companion operator paper anchors its Hilbert--P\'olya
operator.

All statements of this note are formally verified in
\texttt{golden\_mersenne.lean}, a self-contained Lean~4
file built against Mathlib4 (toolchain~\texttt{v4.28.0})~\cite{mathlib};
no \texttt{sorry} occurs in any proof.  Two external inputs enter as
axioms: the Gelfond--Schneider theorem (not yet available in Mathlib
mainline), which underlies the transcendence of $\lambda_{\log}$, and the
44 Lucas--Lehmer~\cite{lehmer} primalities $M_{p}\;\text{prime}$ for
$61\leq p\leq 136\,279\,841$ (GIMPS~\cite{gimps}).  The relation
$\golden^{2}=\golden+1$ is a theorem, not an axiom.

\section{Statement and proof of the universal identity}
\label{sec:universal}

Throughout this section and the next, $\golden^{x}$ for real~$x$
denotes the real exponential $\exp(x\log\golden)$, the standard
extension of integer exponentiation to~$\mathbb{R}$ via the natural
logarithm of the positive real $\golden>0$.

The architectural pattern of~\eqref{eq:euler}---a transcendental
period mediating between algebraic objects---admits a second
realisation on the multiplicative line~$\mathbb{R}_{>0}$, where the
binary base~$2^{p}$ takes the place of the rotor and the golden base
$\golden^{\sigma}$ replaces the angular coordinate.  The mediator is
the constant
\[
  \lambda_{\log} \;:=\; \log 2/\log\golden
                  \;=\; \log_{\golden} 2 \;\approx\; 1.4404,
\]
transcendental over $\mathbb{Q}$ by Gelfond--Schneider
(\Cref{rem:transcendence} below), playing the structural role
of~$2\pi$.  It is the elementary period along which $\golden^{x}$
scales from~$1$ up to the next integer base~$2$:
\begin{equation}\label{eq:golden-turn}
  \golden^{\lambda_{\log}} \;=\; 2,
\end{equation}
the \emph{golden turn} (clause~(iv) of~\Cref{thm:coincidence-intro}).
The multiplicative line~$\mathbb{R}_{>0}$ admits a second canonical
scaling, $\golden^{\log_{\golden} 3} = 3$, and the composition of the
two turns is the central identity of this section:

\begin{Definition}[Golden index]\label{def:sigma}
$\sigma\colon\mathbb{R}\to\mathbb{R},\;\sigma(p):=p\,\lambda_{\log}-\log_{\golden}3.$
\end{Definition}

\begin{Theorem}[Universal identity]\label{thm:universal}
For every $p\in\mathbb{R}$, $\quad 3\,\golden^{\sigma(p)} \;=\; 2^{p}$.
\end{Theorem}

\noindent This is clause~(v) of~\Cref{thm:coincidence-intro}.  The
multiplier~$3$ is not free: \Cref{thm:three-main} below
(clause~(vii)) shows that it coincides with the first Mersenne prime $M_{2}=2^{2}-1$ and with the cardinality of the admissible residue set
of~\Cref{thm:concentration-main}, and no other positive real admits
this simultaneous identification.  The proof of~\Cref{thm:universal}
is established below in three steps:

\begin{Lemma}[Binary turn]\label{lem:binary}
$\golden^{\lambda_{\log}} = 2$.
\end{Lemma}

\begin{proof}
$\golden^{\lambda_{\log}}=\exp(\lambda_{\log}\log\golden)=\exp(\log 2)=2$.
Lean: \lean{mersenne\_bridge\_via\_lambda}.
\end{proof}

\begin{Lemma}[Ternary turn]\label{lem:ternary}
$\golden^{\log_{\golden} 3} = 3$.
\end{Lemma}

\begin{proof}
Identical with $3$ in place of~$2$.  Lean: \lean{phi\_pow\_log3}.
\end{proof}

\begin{proof}[Proof of \Cref{thm:universal}]
By definition $\sigma(p)=p\lambda_{\log}-\log_{\golden}3$, hence
\begin{align*}
  3\,\golden^{\sigma(p)}
  &= 3\,\golden^{p\lambda_{\log}}\,\golden^{-\log_{\golden}3}
  &&\text{(real-exponent subtraction)}\\
  &= 3\,\bigl(\golden^{\lambda_{\log}}\bigr)^{p}\,
       \bigl(\golden^{\log_{\golden}3}\bigr)^{-1}
  &&\text{($\golden^{ab}=(\golden^{a})^{b}$ for $\golden>0$)}\\
  &= 3\cdot 2^{p}\cdot 3^{-1}
  &&\text{(\Cref{lem:binary,lem:ternary})}\\
  &= 2^{p}. \qedhere
\end{align*}
Lean: \lean{golden\_tower\_bridge}.
\end{proof}

\begin{Theorem}[Uniqueness of the index]\label{thm:unique}
For each $p\in\mathbb{R}$, $\sigma(p)$ is the only real~$\sigma$ such that
$3\golden^{\sigma}=2^{p}$.
\end{Theorem}

\begin{proof}
If $3\golden^{\sigma}=2^{p}$ and $3\golden^{\sigma(p)}=2^{p}$ then
$\golden^{\sigma-\sigma(p)}=1$, hence
$(\sigma-\sigma(p))\log\golden=0$, and since $\log\golden>0$ (as
$\golden>1$), $\sigma=\sigma(p)$.  Lean:
\lean{sigma\_mersenne\_unique}.
\end{proof}

\begin{Theorem}[Direct golden ratio between Mersenne]\label{thm:ratio-main}
For every pair $(p,q)\in\mathbb{R}^{2}$,
\[
  \frac{2^{q}}{2^{p}} \;=\; \golden^{(q-p)\lambda_{\log}}.
\]
Specialised to the GIMPS list, for every pair
$(p_{i},p_{j})\in\{p_{1},\ldots,p_{52}\}^{2}$,
\[
  \frac{M_{p_{j}}+1}{M_{p_{i}}+1} \;=\; \golden^{(p_{j}-p_{i})\lambda_{\log}}.
\]
\end{Theorem}

\begin{proof}
By \Cref{lem:binary} and the identity $2^{q}/2^{p}=2^{q-p}$,
\[
  \frac{2^{q}}{2^{p}}
  = 2^{q-p}
  = \bigl(\golden^{\lambda_{\log}}\bigr)^{q-p}
  = \golden^{(q-p)\lambda_{\log}}.
\]
Specialising to GIMPS pairs is immediate.  Lean:
\lean{mersenne\_ratio\_golden} and
\lean{mersenne\_ratio\_at\_known\_pairs}.
\end{proof}

The substance of \Cref{thm:ratio-main} is that the entire orbit
$\{2^{p_{i}}\}_{i=1}^{52}$ fits into a single integer lattice
$\golden^{\mathbb{Z}\lambda_{\log}}$ of the golden base, indexed by the
gaps $p_{j}-p_{i}\in\mathbb{Z}$.  No analogous statement holds for
arbitrary $\{r_{i}\}\subset\mathbb{R}_{>0}$ in place of $\{p_{i}\}$.

\begin{Remark}[The golden lattice in binary]\label{rem:binary}
The ratios of~\Cref{thm:ratio-main} have a transparent binary reading.  In
base~$2$, the Mersenne number $\Mp=2^{p}-1$ is a string of $p$ ones
(\lean{mersenne\_digits\_ones}), while $2^{p}=\Mp+1$ is a single bit---a $1$
followed by $p$ zeros (\lean{two\_pow\_digits\_single\_bit})---the increment
being the carry cascade that turns $\underbrace{1\cdots1}_{p}$ into
$1\underbrace{0\cdots0}_{p}$ (\lean{mersenne\_succ\_eq\_two\_pow}).  The
golden lattice $\{2^{p}\}$ is therefore exactly the set of single-bit
integers, indexed by bit position, and the multiplicative step
$\golden^{\lambda_{\log}}=2$ is the binary left-shift
(\lean{golden\_step\_is\_bit\_shift}).  The golden ratio
$2^{p_{j}}/2^{p_{i}}=\golden^{(p_{j}-p_{i})\lambda_{\log}}$ of
\Cref{thm:ratio-main} is then the shift of the single bit from position
$p_{i}$ to position $p_{j}$.
\end{Remark}

\begin{Corollary}[52-instance specialisation]\label{cor:51}
For each $p\in\{p_{1},\ldots,p_{52}\}$, the unique real~$\sigma$ with
$3\golden^{\sigma}=2^{p}$ is $\sigma=\sigma(p_{i})$; moreover the indices
satisfy $\sigma(p_{1})<\cdots<\sigma(p_{52})$.  Lean:
\lean{mersenne\_phi\_correspondence\_52} and
\lean{sigma\_strictly\_increasing}.
\end{Corollary}

\begin{Remark}[Transcendence of $\lambda_{\log}$]
\label{rem:transcendence}
By Gelfond--Schneider (Gelfond, 1934~\cite{gelfond}; Schneider,
1935~\cite{schneider}), since $\golden$ is algebraic of degree~$2$ over
$\mathbb{Q}$ and $\golden^{\lambda_{\log}}=2$ is algebraic of degree~$1$,
the exponent $\lambda_{\log}$ is transcendental over~$\mathbb{Q}$.  The
identity~\eqref{eq:golden-turn} thus realises an exact algebraic relation
between two algebraic numbers mediated by a transcendental exponent.  This transcendence is formally derived in
Lean~4 (\lean{lambda\_log\_transcendental}), with the Gelfond--Schneider
theorem itself---recently formalised by Karatarakis and Wiedijk~\cite{karatarakis-wiedijk}---taken as the single external
input.
\end{Remark}

\begin{Corollary}[Inverse formula]\label{cor:inverse}
Equivalently, $\sigma(p)=\log_{\golden}(2^{p}/3)$ for every $p\in\mathbb{R}$.
Lean: \lean{sigma\_mersenne\_eq\_logb}.
\end{Corollary}

\begin{Remark}[The index as a golden logarithm]\label{rem:index-logarithm}
The inverse form exhibits the golden index $\sigma(p)$ as a logarithm in
base $\golden$ of an algebraic number.  This is the source of its
transcendence: for every integer $p$, $\golden$ is algebraic of
degree~$2$ and $2^{p}/3$ is rational, so by the same Gelfond--Schneider
argument as for $\lambda_{\log}$ (\Cref{rem:transcendence}),
$\sigma(p)=\log_{\golden}(2^{p}/3)$ is transcendental.  The period
$\lambda_{\log}=\log_{\golden}2$ of \Cref{rem:transcendence} is the
particular golden logarithm against which the whole family $\sigma(p)$ is
measured.
\end{Remark}

\begin{Remark}[Two transcendence theorems, same architectural role]
\label{rem:two-transcendences}
The two periods $2\pi$ of~\eqref{eq:euler} and $\lambda_{\log}$
of~\eqref{eq:golden-turn} are both transcendental but for distinct
reasons: $\pi$ by Lindemann (1882, via the algebraicity
of~$e^{i\pi}=-1$), $\lambda_{\log}$ by Gelfond--Schneider (1934, via
the algebraicity of $\golden^{\lambda_{\log}}=2$).  Two distinct
transcendence theorems play the same structural role in two
realisations of the same architecture: a transcendental period
mediates between algebraic objects in each realisation.
\end{Remark}

\begin{Corollary}[Transcendental injectivity of the golden index]
\label{cor:linear-independence}
For every nonzero integer $p$, $\sigma(p)$ is transcendental
over~$\mathbb{Q}$.  In particular, the 52 values
$\{\sigma(p_{1}),\ldots,\sigma(p_{52})\}\subset\mathbb{R}$ are 52 pairwise
distinct transcendental numbers, occupying 52 distinct points of the
affine $\mathbb{Z}$-lattice $\mathbb{Z}\lambda_{\log}-\log_{\golden}3\subset\mathbb{R}$.
\end{Corollary}

\begin{proof}
By \Cref{cor:inverse}, $\sigma(p)=\log_{\golden}(2^{p}/3)$.  For nonzero
integer~$p$, $2^{p}/3$ is a nonzero rational, and $\golden$ is algebraic of
degree~$2$ with $\golden\notin\{0,1\}$.  By Gelfond--Schneider,
$\golden^{s}$ is transcendental for every algebraic irrational~$s$; since
$\golden^{\sigma(p)}=2^{p}/3$ is rational, $\sigma(p)$ cannot be algebraic
irrational, so it is either rational or transcendental.  It is not
rational: if $\sigma(p)=m/n$ with $m\neq 0$, then
$\golden^{m}=(2^{p}/3)^{n}\in\mathbb{Q}$, but
$\golden^{m}=(L_{m}+F_{m}\sqrt5)/2$ with $F_{m}\neq 0$ for $m\neq 0$, so
$\golden^{m}\notin\mathbb{Q}$; and $m=0$ would give $2^{p}=3$, impossible
for integer~$p$.  Hence $\sigma(p)$ is transcendental.

Pairwise distinctness of $\sigma(p_{1}),\ldots,\sigma(p_{52})$ follows
from the strict monotonicity of~\Cref{cor:51} (a consequence of
$\lambda_{\log}>0$).
\end{proof}

This is the quantitative-arithmetic content of the lattice embedding:
each of the 52 golden coordinates is itself a transcendental, not just
the lattice as a whole, and no algebraic choice of $\lambda$ could
produce such a configuration.  The Mersenne sublattice
$\golden^{\mathbb{Z}\lambda_{\log}-\log_{\golden}3}\subset\mathbb{R}_{>0}$
is genuinely a lattice of $\mathbb{Z}$-rank equal to the number of
distinct Mersenne exponents, with all 52 points transcendental.

\begin{Remark}[Numerical versus algebraic]
A numerical version stating $|3\golden^{\sigma(p)}-2^{p}|<10^{-14}$ over
the 52 known exponents is verifiable in double-precision arithmetic.
\Cref{thm:universal} is strictly stronger: it is exact, holds for all
real~$p$, and is independent of floating-point resolution.
\end{Remark}

\section{The factor 3: the triple coincidence}
\label{sec:triple}

\begin{Theorem}[Modular concentration]\label{thm:concentration-main}
For every prime~$p$, $\;\Mp\bmod 20 \in \{3,7,11\}$, a subset of
cardinality~$3=M_{2}$ of the eight residue classes
$\Ztwenty=\{1,3,7,9,11,13,17,19\}$.  Lean:
\lean{mersenne\_concentration\_general}.
\end{Theorem}

\begin{proof}
For $p\geq 5$ prime, $p$ is odd and $p\bmod 4\in\{1,3\}$.  Two
elementary facts in $\mathbb{Z}/20\mathbb{Z}$: $2^{4}\equiv 16$ and
$16^{k}\equiv 16$ for all $k\geq 1$.  Writing $p=4\lfloor p/4\rfloor+r$
with $r\in\{1,3\}$,
$
  2^{p} = 16^{\lfloor p/4\rfloor}\cdot 2^{r} \equiv 16\cdot 2^{r}
  \pmod{20},
$
so $2^{p}\bmod 20\in\{12,8\}$ and $\Mp\bmod 20\in\{11,7\}$.  The cases
$p=2$, $p=3$ give $M_{2}=3$, $M_{3}=7$ directly; combining,
$\Mp\bmod 20\in\{3,7,11\}$.
\end{proof}

\begin{Remark}[Relation to prior
literature]\label{rem:acevedo}
The concentration of~\Cref{thm:concentration-main} is a corollary of
the stronger classification of Acevedo
Agudelo~\cite{acevedo-2020}:
$\Mp\bmod 360\in\{3,7,31,127,247,271\}$ for every prime~$p$, which
reduces modulo~$20$ to the present statement. Partial precedents in
the OEIS comments for sequence A000668 record weaker congruences
($\Mp\equiv 1\pmod{6}$, $\Mp\equiv 7\pmod{12}$,
$\Mp\equiv 7\pmod{24}$), none of which alone implies the mod-$20$
concentration. The role of~\Cref{thm:concentration-main} in the
present note is as clause~(vi)
of~\Cref{thm:euler-connector}, conjoined with clauses~(vii) and~(xi).
\end{Remark}

The restriction of all $\Mp$ to $3$ of the $8$ residue classes coprime
with~$20$ is not visible from the definition of the Mersenne sequence;
it is a phenomenon of the modular doubling pattern $2^{p}\bmod 20$
produced by the proof above.  The numerator and denominator of the
ratio $3/8=M_{2}/|\Ztwenty|$ that records this concentration
quantitatively are themselves clauses~(vii) and~(ix)
of~\Cref{thm:euler-connector}.

\begin{Theorem}[The factor~$3$ is structurally $M_{2}$]\label{thm:three-main}
The natural number~$3$ admits simultaneously the three identifications
\[
  3 \;=\; M_{2}=2^{2}-1 \;=\; \bigl|\{3,7,11\}\bigr|.
\]
\end{Theorem}

\begin{proof}
The chain $3 = M_{2} = 2^{2}-1 = |\{3,7,11\}|$ comprises three
identifications.  First, $M_{2}=2^{2}-1$ by definition of the Mersenne
sequence, and $2^{2}-1=3$ by computation.  Second,
$|\{3,7,11\}|=3$ by counting.  Third, that $\{3,7,11\}$ is exactly
the set of residue classes containing $\Mp\bmod 20$ for some
prime~$p$ is~\Cref{thm:concentration-main}.  Lean:
\lean{three\_eq\_M2\_eq\_admissible\_card} and
\lean{admissible\_classes\_card\_eq\_M2}.
\end{proof}

The asymmetry between $2$ and $3$ in \Cref{thm:universal} is only
apparent: $2$ enters as the base of $2^{p}$ (the substrate), while $3$
enters as the multiplier.  The triple identification of
\Cref{thm:three-main} shows that the choice of multiplier is not free.
Replacing $3$ by any other positive real~$c$ in the identity produces
$c\,\golden^{p\lambda_{\log}-\log_{\golden}c}=2^{p}$, but no other value
of~$c$ admits the simultaneous identification with $M_{2}$, with the
cardinality of the admissible residue classes of~$\Mp\bmod 20$, and with
the smallest Mersenne prime occurring as the residue
$M_{2}\equiv 3\pmod{20}$ in the admissible set $\{3,7,11\}$.  The value
$c=3$ is determined by the arithmetic of Mersenne in~$\Ztwenty$, not by
the algebra of the identity.

\begin{Remark}[Pentagonal--cyclotomic underpinning]\label{rem:pentagon}
The golden ratio admits the pentagonal identity
\begin{equation}\label{eq:pentagon}
  \golden \;=\; 2\cos(\pi/5) \;=\; \zeta_{10}+\zeta_{10}^{-1},
  \qquad \zeta_{10} := e^{i\pi/5},
\end{equation}
where $\zeta_{10}^{10}=1$.  The first equality
$\golden=2\cos(\pi/5)$ derives from the Chebyshev relation
$T_{5}(\cos\theta)=\cos(5\theta)$ at $\theta=\pi/5$, which gives
$4\cos^{2}(\pi/5)-2\cos(\pi/5)-1=0$ (after factoring out the root
$\cos(\pi/5)\neq -1$), and selecting the positive root yields
$\cos(\pi/5)=\golden/2$.  The second equality
$2\cos(\pi/5)=\zeta_{10}+\zeta_{10}^{-1}$ is the real part of Euler's
identity~\eqref{eq:euler} at $\theta=\pi/5$:
$e^{i\pi/5}+e^{-i\pi/5}=2\cos(\pi/5)$.

These two equalities, taken jointly, register four of the twelve
clauses of~\Cref{thm:euler-connector}.  Clause~(iii)
$\golden=2\cos(\pi/5)$ is the real-trigonometric embodiment of the
algebraic object: the same $\golden$ defined by $\golden^{2}=\golden+1$,
which carries no intrinsic notion of rotation or periodicity, here
appears inscribed in~$\mathbb{R}$ through cosine.  Clause~(i)
$\golden_\mathbb{C}=\zeta_{10}+\zeta_{10}^{-1}$ is the
cyclotomic-complex embodiment: the same algebraic object now appears as
a concrete point of Euler's structure---the sum of two conjugate
vertices of the regular decagon $\{\zeta_{10}^{k}\}_{k=0}^{9}$
in~$S^{1}$.  Clause~(ii) $\zeta_{10}^{10}=1$ places $\zeta_{10}$ not
as a free point of $S^{1}$ but as a generator of the cyclic group of
order~$10$; without it $\zeta_{10}$ would admit no Galois structure.
Equivalently, $\golden$ inhabits the real maximal subfield
$\mathbb{Q}(\golden)\subset\mathbb{Q}(\zeta_{20})$ of the cyclotomic
field whose Galois group is
\[
  \mathrm{Gal}(\mathbb{Q}(\zeta_{20})/\mathbb{Q}) \;\cong\;
  \Ztwenty,
\]
an abelian group of order $\phi(20)=8$; this is clause~(ix), the
cardinality of the Galois receptacle within which the Mersenne
residues of~\Cref{thm:concentration-main} are hosted.  The four
clauses (i),~(ii),~(iii),~(ix) present the same algebraic
object~$\golden$ in four canonical inscriptions inside the structure
of~\eqref{eq:euler}.  The modular concentration of
\Cref{thm:concentration-main} into~$3$ classes of~$\Ztwenty$, together
with the identification $3=M_{2}$ of~\Cref{thm:three-main}, are the
arithmetic shadows of this geometric inscription.  Formal identities:
\lean{phi\_eq\_two\_cos\_pi\_fifth}, \lean{phi\_eq\_zeta10\_sum},
\lean{zeta10\_pow\_ten\_eq\_one}.
\end{Remark}

\begin{Theorem}[Simultaneous identities theorem for
$\golden$]\label{thm:euler-connector}
A single $\golden$ satisfies simultaneously the following twelve
identities, inclusions, and identifications in their respective
canonical structures:
\begin{enumerate}[topsep=0pt,itemsep=0.1em,leftmargin=2.3em]
  \item[\textup{(i)}] \emph{Cyclotomic inscription in $\mathbb{C}$:}
    $\golden_\mathbb{C} = \zeta_{10}+\zeta_{10}^{-1}$, with
    $\zeta_{10}:=e^{i\pi/5}$ (\Cref{rem:pentagon}; Lean:
    \lean{phi\_eq\_zeta10\_sum}).
  \item[\textup{(ii)}] \emph{Tenth-root closure:} $\zeta_{10}^{10}=1$
    (Lean: \lean{zeta10\_pow\_ten\_eq\_one}).
  \item[\textup{(iii)}] \emph{Pentagonal trigonometric identity in
    $\mathbb{R}$:} $\golden = 2\cos(\pi/5)$ (\Cref{rem:pentagon}; Lean:
    \lean{phi\_eq\_two\_cos\_pi\_fifth}).
  \item[\textup{(iv)}] \emph{Logarithmic binary bridge:}
    $\golden^{\lambda_{\log}} = 2$ (\Cref{lem:binary}; Lean:
    \lean{mersenne\_bridge\_via\_lambda}).
  \item[\textup{(v)}] \emph{Universal real-multiplicative identity:}
    for every $p\in\mathbb{R}$, $3\,\golden^{\sigma(p)} = 2^{p}$
    (\Cref{thm:universal}; Lean: \lean{golden\_tower\_bridge}).
  \item[\textup{(vi)}] \emph{Modular concentration:} for every prime
    $p$, $\Mp\bmod 20\in\{3,7,11\}$ (\Cref{thm:concentration-main};
    Lean: \lean{mersenne\_concentration\_general}).
  \item[\textup{(vii)}] \emph{Triple ternary identification:}
    $3 = M_{2} = |\{3,7,11\}|$ (\Cref{thm:three-main}; Lean:
    \lean{three\_eq\_M2\_eq\_admissible\_card}).
  \item[\textup{(viii)}] \emph{Real-complex cross identification:}
    $\iota(\golden)=\golden_\mathbb{C}$.
  \item[\textup{(ix)}] \emph{Galois order:}
    $|\Ztwenty|=8$ (Lean: \lean{ZtwentyStar\_card}).
  \item[\textup{(x)}] \emph{Galois sub-object inclusion:}
    $\{3,7,11\}\subseteq\Ztwenty$.
  \item[\textup{(xi)}] \emph{Quantitative ratio:}
    \[
      \frac{|\{3,7,11\}|}{|\Ztwenty|}
      \;=\; \frac{M_{2}}{8} \;=\; \frac{3}{8},
    \]
    equivalently $|\{3,7,11\}|\cdot 8 = M_{2}\cdot|\Ztwenty|$.
  \item[\textup{(xii)}] \emph{Rotor discretisation:} for every
    $k\in\mathbb{N}$, $(\zeta_{10}^{k})^{10}=1$.
\end{enumerate}
Formal synthesis: \lean{phi\_euler\_mersenne\_connector}.
\end{Theorem}

Clauses (viii), (x), (xi), (xii) have short proofs but each registers
a structural knotting that the theorem requires.

\textbf{Clause~(viii)} (real-complex cross identification) is the
identification that permits speaking of \emph{the same}~$\golden$
across clause~(iii) (trigonometric in~$\mathbb{R}$) and clause~(i)
(cyclotomic in~$\mathbb{C}$).  Without it, the statement would not be
about a single algebraic object but about distinct copies
of~$\golden$ in disjoint structures.  The proof is the definitional
behaviour of the canonical embedding $\iota$ on the algebraic
point~$\golden$.

\textbf{Clause~(x)} ($\{3,7,11\}\subseteq\Ztwenty$) registers that
the three admissible Mersenne residues of clause~(vi) are not
arbitrary integers but elements of the Galois group of clause~(ix):
the arithmetic locus of Mersenne primes is inscribed inside the
cyclotomic Galois receptacle.  The proof is that $3$, $7$, $11$ are
each coprime with~$20$.

\textbf{Clause~(xi)} ($|\{3,7,11\}|/|\Ztwenty|=3/8$) is the
quantitative knotting: the first Mersenne prime~$M_{2}$ is the
numerator of the fraction of Galois classes that host all Mersenne
arithmetic.  The proof is the direct numerical consequence of
clauses~(vii) ($M_{2}=3$) and~(ix) ($|\Ztwenty|=8$); the substance is
that the simultaneous holding of those two clauses under the
same~$\golden$ is what produces the ratio $3/8$ as structural and not
as numerical coincidence.

\textbf{Clause~(xii)} ($(\zeta_{10}^{k})^{10}=1$ for all
$k\in\mathbb{N}$) registers the canonical discretisation of the
continuous rotation $\theta\mapsto e^{i\theta}$ of~\eqref{eq:euler}:
evaluated at $\theta=2\pi k/10$ it produces a cyclic group of
order~$10$.  Clause~(xii) is the first knotting in the statement
between continuous and discrete: without it, the decagonal vertices
$\{\zeta_{10}^{k}\}_{k=0}^{9}$ are a set of ten complex numbers; with
it, they are a finite cyclic group on which the Galois structure
of~(ix) acts.  The proof is
$(\zeta_{10}^{k})^{10}=(\zeta_{10}^{10})^{k}=1^{k}=1$ by clause~(ii).

All four are subsumed in the formal
synthesis~\lean{phi\_euler\_mersenne\_connector}.

\begin{Remark}[Reading of \Cref{thm:euler-connector}]\label{rem:euler-reading}
No proper sub-collection of the twelve clauses constitutes the content
of \Cref{thm:euler-connector}; the content is the simultaneous holding
of all twelve under a single $\golden$.  The clauses inhabit canonical
structures of distinct kinds: complex-rotational and cyclotomic
(i,~ii,~xii), real-trigonometric (iii), logarithmic and
real-multiplicative (iv,~v), arithmetic-modular over $\mathbb{Z}/20\mathbb{Z}$
(vi,~vii), cross-domain real-to-complex (viii), Galois-theoretic of the
cyclotomic decagonal extension $\mathbb{Q}(\zeta_{20})/\mathbb{Q}$
(ix,~x,~xi).  The continuous rotational symmetry of
Euler's identity~\eqref{eq:euler}, with period $2\pi$ in
$S^{1}\subset\mathbb{C}$, discretises canonically into the ten
vertices of the regular decagon
$\{\zeta_{10}^{k}\}_{k=0}^{9}$ (clause~xii), of which $\golden$ is the
sum of two conjugate vertices (clauses i,~iii); the Galois group of
the larger cyclotomic extension is~$\Ztwenty$ of order~$8$
(clause~ix); the Mersenne primes $\Mp$ occupy $M_{2}=3$ of those $8$
classes (clauses vi,~vii,~x,~xi); and the same $\golden$ mediates, via
the transcendental exponent~$\lambda_{\log}$ (\Cref{rem:transcendence}),
the binary doubling $\golden^{\lambda_{\log}}=2$ and the universal
identity $3\golden^{\sigma(p)}=2^{p}$ (clauses iv,~v).  The continuous
rotor, the discrete arithmetic of Mersenne, the cyclotomic geometry of
$\mathbb{Q}(\zeta_{20})$, and the real-multiplicative line
$\mathbb{R}_{>0}$ coincide simultaneously under~$\golden$ across the
twelve clauses of~\Cref{thm:euler-connector}.  The quantitative ratio
$M_{2}/|\Ztwenty|=3/8$ of clause~(xi) records that the first Mersenne
prime is the numerator of the fraction of Galois classes that host
the Mersenne arithmetic.

The arithmetic of the integers, in this reading, appears as a discrete
shadow of the rotational symmetry of~\eqref{eq:euler}: the chain
rotation~$\to$ polygonal discretisation~$\to$ Galois action~$\to$
residue classification~$\to$ Mersenne concentration formalises a
geometric connection with the tools of Gelfond--Schneider and
\texttt{AdjoinRoot}.
\end{Remark}

\subsection{The identity, characterized}
\label{ssec:characterization}

Read as an equality, the universal identity of~\Cref{thm:universal} is the
evaluation of a definition: $\sigma(p):=p\,\lambda_{\log}-\log_{\golden}3$ is
chosen precisely so that $3\,\golden^{\sigma(p)}=2^{p}$.  What is \emph{not}
definitional is the statement gathered here: that the constituents of the
equation---the factor $3$, the base $\golden$, the slope
$\lambda_{\log}$---are each pinned by a theorem \emph{external} to the
equation, and that these separately fixed determinations are mutually
consistent, with the golden coordinate uniquely determined.

\begin{Theorem}[The identity as a consistency theorem]\label{thm:characterization}
The equation $3\,\golden^{\sigma}=2^{p}$ in the unknown $\sigma\in\mathbb{R}$
carries four determinations, each fixed independently of it, and they hold
simultaneously:
\begin{enumerate}
\item[\textup{(i)}] \emph{Existence.} $\sigma(p)$ solves it for every real
  $p$ (\Cref{thm:universal}).
\item[\textup{(ii)}] \emph{Uniqueness.} $\sigma(p)$ is its only real
  solution (\Cref{thm:unique}).
\item[\textup{(iii)}] \emph{Transcendence.} the slope $\lambda_{\log}$ is
  transcendental over $\mathbb{Q}$ (\Cref{rem:transcendence}).
\item[\textup{(iv)}] \emph{Forced factor.} $3=M_{2}=|\{3,7,11\}|$, the number
  of residue classes the Mersenne primes occupy modulo~$20$
  (\Cref{thm:concentration-main}), independently of the equation.
\end{enumerate}
\emph{Lean:} \lean{golden\_identity\_characterization}.
\end{Theorem}

\begin{Remark}[What the theorem certifies]\label{rem:characterization}
The central equation~(i) is not proved against itself---that would be
circular; what is proved is that it exists, is unique, and coincides with
the external determinations~(iii)--(iv).  The uniqueness~(ii) certifies that
the golden coordinate $\sigma(p)$ is well defined; the evidential weight
lies in~(iii) and~(iv), where the slope and the factor are pinned by
theorems that never mention the equation.  In this reading the identity is
neither an axiom nor a bare definition, but the single point at which
independently established facts about $\golden$ agree.  Formally, $\sigma$
is the unique function with $3\,\golden^{\sigma(p)}=2^{p}$ for all~$p$
(\lean{golden\_identity\_is\_consistent\_overdetermined}).
\end{Remark}


\section{Fibonacci splitting of the Mersenne primes}
\label{sec:fibonacci}

The modular concentration of~\Cref{thm:concentration-main} records
\emph{where} the Mersenne primes fall in $\Ztwenty$: their residues lie
in $\{3,7,11\}$.  We now show \emph{how} they fall---how each prime
$\Mp$ splits in the field $\mathbb{Q}(\sqrt5)$ of the golden ratio.  The
answer is governed by a single classical congruence (Lucas, 1878) and,
composed with the concentration, refines the three admissible classes
into a split class and two inert classes.  This is the arithmetic
content behind the identification $3=M_{2}=|\{3,7,11\}|$ of
clause~(vii): the multiplier~$3$ is the number of classes the Mersenne
primes occupy, and the golden field determines the splitting type in
each.

\subsection{Fibonacci prerequisites}
\label{ssec:fib-prereq}

The Fibonacci sequence is $F_{0}=0$, $F_{1}=1$, $F_{n+2}=F_{n+1}+F_{n}$.
In any commutative ring containing an element~$\alpha$ with
$\alpha^{2}=\alpha+1$, induction on~$n$ gives the Binet-type identity
\begin{equation}\label{eq:binet}
  \alpha^{\,n+1} = F_{n+1}\,\alpha + F_{n}.
\end{equation}
The Galois conjugate of~$\golden$ over~$\mathbb{Q}$ is
$\overline{\golden}=1-\golden=-1/\golden$, with trace
$\golden+\overline{\golden}=1$ and norm
$\golden\,\overline{\golden}=-1$; both roots of $X^{2}-X-1$.  The
discriminant of $\mathbb{Q}(\sqrt5)/\mathbb{Q}$ is
$(\golden-\overline{\golden})^{2}=5$, so~$5$ is the unique ramified
prime.  For an odd prime $q\neq5$, the splitting of~$q$ in
$\mathbb{Q}(\sqrt5)$ is governed by the Legendre symbol $(q/5)$:
$(q/5)=+1$ (split) iff $q\bmod 20\in\{1,9,11,19\}$, and $(q/5)=-1$
(inert) iff $q\bmod 20\in\{3,7,13,17\}$.  Among the admissible Mersenne
classes~$\{3,7,11\}$, exactly one---the class~$11$---is split.
Lean: \lean{fib\_pow\_aux}, \lean{chi5\_values}.

\subsection{The Fibonacci splitting congruence}
\label{ssec:fib-congruence}

\begin{Theorem}[Fibonacci splitting; Lucas 1878]\label{thm:fib-main}
For every prime $q\neq2,5$,
\[
  F_{q}\;\equiv\;(q/5)\pmod{q},
\]
with value $+1$ in the split case and $-1$ in the inert case.
\emph{Lean:} \lean{fibonacci\_splitting\_general}.
\end{Theorem}

\begin{Lemma}[Split case]\label{lem:fib-split}
Let $q\neq2,5$ be prime with $(q/5)=+1$.  Then $X^{2}-X-1$ has a root
$\alpha\in\mathbb{F}_{q}$, and $F_{q}\equiv 1\pmod{q}$.
\emph{Lean:} \lean{fib\_split\_case}.
\end{Lemma}

\begin{proof}
Since $(q/5)=+1$, there is $\beta\in\mathbb{F}_{q}$ with $\beta^{2}=5$;
put $\alpha:=(1+\beta)/2$ (defined as $q\neq2$).  Then
$\alpha^{2}=(6+2\beta)/4=\alpha+1$.  By Fermat, $\alpha^{q}=\alpha$, and
by~\eqref{eq:binet}, $\alpha^{q}=F_{q}\alpha+F_{q-1}$; the same holds for
the conjugate root $\overline{\alpha}=1-\alpha$.  Subtracting the two
relations gives $F_{q}(\alpha-\overline{\alpha})\equiv
\alpha-\overline{\alpha}$; since
$(\alpha-\overline{\alpha})^{2}=5\not\equiv0$ (as $q\neq5$) the factor is
invertible, so $F_{q}\equiv1\pmod{q}$.
\end{proof}

\paragraph{Worked example ($q=11$, split).}
Since $11\equiv1\pmod5$, $(11/5)=+1$; in $\mathbb{F}_{11}$ the roots of
$X^{2}=X+1$ are $\alpha=8$, $\overline{\alpha}=4$ (as
$8^{2}=64\equiv9=8+1$).  And $F_{11}=89=11\cdot8+1$, so
$F_{11}\equiv1\pmod{11}$.

\begin{Lemma}[Inert case]\label{lem:fib-inert}
Let $q\neq2,5$ be prime with $(q/5)=-1$.  Then $X^{2}-X-1$ is
irreducible over~$\mathbb{F}_{q}$, the Frobenius $x\mapsto x^{q}$ of
$\mathbb{F}_{q^{2}}=\mathbb{F}_{q}[X]/(X^{2}-X-1)$ exchanges the two
roots $\alpha,\overline{\alpha}=1-\alpha$, and $F_{q}\equiv-1\pmod{q}$.
\emph{Lean:} \lean{fib\_inert\_case}.
\end{Lemma}

\begin{proof}
By quadratic reciprocity, $(5/q)=(q/5)=-1$, so~$5$ is a non-residue
and $X^{2}-X-1$ has no root in~$\mathbb{F}_{q}$; being quadratic it is
irreducible, and the quotient
$\mathbb{F}_{q^{2}}=\mathbb{F}_{q}[X]/(X^{2}-X-1)$ is a field of
order~$q^{2}$ with basis $\{1,\alpha\}$, where $\alpha$ is the image
of~$X$.  The Frobenius $\mathrm{Frob}_{q}\colon x\mapsto x^{q}$ is an
$\mathbb{F}_{q}$-automorphism (additivity is the freshman's dream in
characteristic~$q$; fixity of~$\mathbb{F}_{q}$ is Fermat), it permutes
the roots of $X^{2}-X-1$, and it is not the identity (else all of
$\mathbb{F}_{q^{2}}$ would lie in~$\mathbb{F}_{q}$); hence
$\alpha^{q}=\overline{\alpha}=1-\alpha$.  By~\eqref{eq:binet},
$\alpha^{q}=F_{q}\alpha+F_{q-1}$, so $(F_{q}+1)\alpha+(F_{q-1}-1)=0$ in
$\mathbb{F}_{q^{2}}$; as $\{1,\alpha\}$ is a basis, both coefficients
vanish, giving $F_{q}\equiv-1\pmod{q}$.
\end{proof}

\begin{Remark}[Lucas--Lehmer companion]\label{rem:fib-companion}
The inert case also yields $F_{q-1}\equiv1$, whence
$F_{q+1}=F_{q}+F_{q-1}\equiv0\pmod{q}$: the classical Fermat--Lucas
condition $q\mid F_{q+1}$ for inert primes, which underlies the index of
appearance in the Lucas--Lehmer primality test.
\end{Remark}

\paragraph{Worked example ($q=7$, inert).}
Since $7\equiv2\pmod5$ is a non-residue, $(7/5)=-1$ and $X^{2}-X-1$ is
irreducible over~$\mathbb{F}_{7}$.  Here $F_{7}=13\equiv-1\pmod{7}$, and
$F_{8}=21\equiv0\pmod7$, consistent with~\Cref{rem:fib-companion}.

\subsection{Splitting type of the Mersenne primes}
\label{ssec:fib-mersenne}

\begin{Theorem}[Residue class determines splitting type]\label{thm:residue-split}
Let~$p$ be a prime with $\Mp$ also prime, $\Mp\neq2,5$.  Then
$\Mp\bmod20\in\{3,7,11\}$, and
\[
  \Mp\equiv11\Rightarrow F_{\Mp}\equiv+1\ (\text{split}),
  \qquad
  \Mp\bmod20\in\{3,7\}\Rightarrow F_{\Mp}\equiv-1\ (\text{inert}),
\]
all congruences modulo~$\Mp$.
\emph{Lean:} \lean{mersenne\_fibonacci\_type}.
\end{Theorem}

\begin{proof}
Since $\Mp=2^{p}-1$ is odd and $2^{p}-1=5$ has no solution, the hypothesis $\Mp\neq2,5$ is automatic for every Mersenne prime.  The concentration is~\Cref{thm:concentration-main} at~$p$.
By~\Cref{thm:fib-main}, $F_{\Mp}\equiv(\Mp/5)\pmod{\Mp}$; and
$(\Mp/5)$ is read from $\Mp\bmod5$: $11\equiv1$ gives $+1$, while
$3\equiv3$ and $7\equiv2$ are non-residues, giving $-1$.
\end{proof}

\begin{table}[ht]\centering\footnotesize
\begin{tabular}{@{}cccccc@{}}
\toprule
$p$ & $\Mp$ & $\Mp\bmod20$ & $(\Mp/5)$ & $F_{\Mp}\bmod\Mp$ & type \\
\midrule
$3$  & $7$          & $7$  & $-1$ & $-1$ & inert \\
$5$  & $31$         & $11$ & $+1$ & $+1$ & split \\
$7$  & $127$        & $7$  & $-1$ & $-1$ & inert \\
$13$ & $8191$       & $11$ & $+1$ & $+1$ & split \\
$17$ & $131071$     & $11$ & $+1$ & $+1$ & split \\
$19$ & $524287$     & $7$  & $-1$ & $-1$ & inert \\
$31$ & $2147483647$ & $7$  & $-1$ & $-1$ & inert \\
\bottomrule
\end{tabular}
\caption{Fibonacci splitting type of $\Mp$ for the first seven Mersenne
primes with $p\geq3$.}\label{tab:mersenne-fib}
\end{table}

Specialised to each of the 52 known Mersenne exponents, this classifies
every known Mersenne prime by its splitting type; the distribution
is~\Cref{tab:splitting}.

\begin{table}[ht]\centering\footnotesize
\begin{tabular}{@{}cccc@{}}
\toprule
$\Mp\bmod20$ & Type in $\mathbb{Q}(\sqrt5)$ & $F_{\Mp}\bmod\Mp$ & Count \\
\midrule
$3$  & inert & $-1$ & $1$  \\
$7$  & inert & $-1$ & $19$ \\
$11$ & split & $+1$ & $32$ \\
\midrule
\multicolumn{3}{@{}r}{\textbf{total}} & \textbf{52} \\
\bottomrule
\end{tabular}
\caption{Splitting type in $\mathbb{Q}(\sqrt5)$ of each of the 52 known
Mersenne primes, by residue class $\Mp\bmod20$.  Verified in Lean by
\texttt{native\_decide} for $p\leq31$ and derived from
\lean{mersenne\_concentration\_general} otherwise.}\label{tab:splitting}
\end{table}

\begin{sloppypar}\small
\noindent\emph{Explicit lists.}  The single exponent in class~$3$ is
$p=2$ (inert).  The 19 exponents with $\Mp\equiv7\pmod{20}$ (inert) are
$3$, $7$, $19$, $31$, $107$, $127$, $607$, $1279$, $2203$, $4423$, $86\,243$, $110\,503$, $216\,091$, $756\,839$, $1\,257\,787$, $20\,996\,011$, $24\,036\,583$, $25\,964\,951$, $37\,156\,667$.
The 32 exponents with $\Mp\equiv11\pmod{20}$ (split) are
$5$, $13$, $17$, $61$, $89$, $521$, $2281$, $3217$, $4253$, $9689$, $9941$, $11\,213$, $19\,937$, $21\,701$, $23\,209$, $44\,497$, $132\,049$, $859\,433$, $1\,398\,269$, $2\,976\,221$, $3\,021\,377$, $6\,972\,593$, $13\,466\,917$, $30\,402\,457$, $32\,582\,657$, $42\,643\,801$, $43\,112\,609$, $57\,885\,161$, $74\,207\,281$, $77\,232\,917$, $82\,589\,933$, $136\,279\,841$.
\end{sloppypar}

\begin{Remark}[Classification under equivalence]\label{rem:classification}
Together, \Cref{thm:concentration-main} and~\Cref{thm:residue-split}
realise a classification of the Mersenne primes within the
pentagonal--cyclotomic geometry of~$\golden$: the multiplier
$3=M_{2}=|\{3,7,11\}|$ counts the admissible classes, and the golden
field partitions them into one split class ($11$) and two inert classes
($3,7$).  The observed split/inert split---$32$ against $20$ over the
current 52---is empirical; under conjectural equidistribution one would
expect a ratio nearer $1:2$, and the deviation over a sample of size~$52$
is not conclusive.
\end{Remark}


\subsection{The splitting type from the exponent}
\label{ssec:fib-exponent}

The residue $\Mp\bmod20$ of~\Cref{thm:residue-split} is itself fixed by the
exponent modulo~$4$, which reduces the whole classification to the
arithmetic of~$p$ alone.

\begin{Corollary}[Splitting type from $p\bmod4$]\label{cor:split-from-p}
Let $\Mp$ be a Mersenne prime.  Its splitting type in $\mathbb{Q}(\sqrt5)$
is determined by $p\bmod4$:
\[
  p\equiv1\!\!\pmod4\ \Rightarrow\ \text{split},
  \qquad
  p\equiv3\!\!\pmod4\ \text{or}\ p=2\ \Rightarrow\ \text{inert}.
\]
  \emph{Lean:} \lean{mersenne\_mod5\_from\_p\_mod4}.
\end{Corollary}

\begin{proof}
The residues $2^{p}\bmod5$ have period~$4$, equal to $2,4,3,1$ for
$p\equiv1,2,3,0\pmod4$.  Hence $\Mp=2^{p}-1$ modulo~$5$ depends only on
$p\bmod4$, taking the values $1,3,2$ for $p\equiv1,2,3\pmod4$, so
$(\Mp/5)=+1,-1,-1$ respectively.  By~\Cref{thm:residue-split} this sign is
the splitting type.  (The case $p\equiv0\pmod4$ gives $5\mid\Mp$ and is
excluded, since $\Mp$ is a prime $\neq5$.)
\end{proof}

Thus, among the admissible classes, the split class~$11$ is exactly
$p\equiv1\pmod4$ and the inert classes $3,7$ are $p=2$ and $p\equiv3\pmod4$.
The finer modulus~$8$ does not separate the type any further: $p\bmod4$ is
the exact controlling residue of the exponent.

\begin{Corollary}[Fibonacci rank characterises the inert primes]\label{prop:rank-inert}
For a Mersenne prime $\Mp\neq5$,
\[
  \Mp\mid F_{\Mp+1}
  \;\iff\;
  \Mp\ \text{is inert in }\mathbb{Q}(\sqrt5)
  \;\iff\;
  p\equiv3\!\!\pmod4\ \text{or}\ p=2.
\]
  \emph{Lean:} \lean{mersenne\_rank\_iff\_inert}.
\end{Corollary}

\begin{proof}
If $\Mp$ is inert then $\Mp\mid F_{\Mp+1}$ by~\Cref{rem:fib-companion}.  If
$\Mp$ is split, the roots $\alpha,\overline{\alpha}$ of $X^{2}-X-1$ lie in
$\mathbb{F}_{\Mp}^{\times}$, so $\alpha^{\Mp-1}=\overline{\alpha}^{\Mp-1}=1$
by Fermat, giving
$F_{\Mp-1}=(\alpha^{\Mp-1}-\overline{\alpha}^{\Mp-1})/(\alpha-\overline{\alpha})\equiv0$;
with $F_{\Mp}\equiv1$ (\Cref{lem:fib-split}) this yields
$F_{\Mp+1}=F_{\Mp}+F_{\Mp-1}\equiv1\not\equiv0\pmod{\Mp}$.  The second
equivalence is~\Cref{cor:split-from-p}.
\end{proof}

The divisibility $\Mp\mid F_{\Mp+1}$---that is, the rank of apparition of
$\Mp$ in the Fibonacci sequence dividing $\Mp+1$---therefore singles out
\emph{exactly} the inert half of the Mersenne primes, the same
rank-of-apparition structure that underlies the Lucas--Lehmer test.

\begin{Corollary}[The joint presentation from the exponent]\label{cor:joint}
For a Mersenne prime $\Mp$, the exponent residue $p\bmod4$ fixes the
triple (residue class modulo~$20$, splitting type in $\mathbb{Q}(\sqrt5)$,
Fibonacci sign $F_{\Mp}\bmod\Mp$) at once:
\[
  \begin{aligned}
    p\equiv1\!\!\pmod4 &\mapsto(11,\ \mathrm{split},\ +1),\\
    p\equiv3\!\!\pmod4 &\mapsto(7,\ \mathrm{inert},\ -1),\\
    p=2 &\mapsto(3,\ \mathrm{inert},\ -1).
  \end{aligned}
\]
Over the 52 known Mersenne primes the triple takes exactly these three
values, one per admissible residue; given $p\bmod4$, no further datum
distinguishes the class, the type, and the sign.
\begin{proof}
The class follows from the modular concentration
(\Cref{thm:concentration-main}) together with $\Mp\bmod5$, which
$p\bmod4$ fixes (\Cref{cor:split-from-p}); the type from the same
residue; and the Fibonacci sign from quadratic reciprocity, since
$X^{2}-X-1$ has a root modulo $\Mp$ exactly when $5$ is a square there,
i.e.\ when $\Mp\equiv\pm1\pmod5$, whence \Cref{lem:fib-split} or
\Cref{lem:fib-inert} applies.
\end{proof}
\emph{Lean:} \lean{mersenne\_presentation\_from\_p\_mod4}.
\end{Corollary}


\section{\texorpdfstring{The reciprocal turn $\golden^{-\lambda_{\log}}=\tfrac12$}{The reciprocal turn}}
\label{sec:spectral}

The transcendental period $\lambda_{\log}=\log_{\golden}2$ has two
directions.  Forward it closes the binary bridge
$\golden^{\lambda_{\log}}=2$ of clause~(iv) (\Cref{lem:binary}); read
backward it closes the reciprocal turn
$\golden^{-\lambda_{\log}}=\tfrac12$.  This value is not a thirteenth
clause but a point already inscribed in the identity, legible at three
of its nodes: it is the reciprocal of clause~(iv), the $p=-1$ node of
the universal identity~(v) (\Cref{thm:universal}), and the
multiplicative-line image of the conjugate root of the rotor
(clause~(i)).  We establish the three, then note where the value leads.

\begin{Lemma}[Reciprocal face of the binary bridge]\label{lem:inverse-turn}
$\golden^{-\lambda_{\log}}=\tfrac12$.
\end{Lemma}

\begin{proof}
By clause~(iv) (\Cref{lem:binary}), $\golden^{\lambda_{\log}}=2$; hence
$\golden^{-\lambda_{\log}}=2^{-1}=\tfrac12$.  In Lean this is the inverse
of \lean{mersenne\_bridge\_via\_lambda}.
\end{proof}

\begin{Corollary}[The reciprocal turn is the universal identity at $p=-1$]\label{cor:reciprocal-node}
The reciprocal turn $\golden^{-\lambda_{\log}}=\tfrac12$ is the value of
the universal identity~(v) (\Cref{thm:universal}) at $p=-1$.
\end{Corollary}

\begin{proof}
The universal identity reads $3\,\golden^{\sigma(p)}=2^{p}$ with
$\sigma(p)=p\,\lambda_{\log}-\log_{\golden}3$.  At $p=-1$ the binary side
is $2^{-1}=\tfrac12$, and the golden side reduces as
\[
  3\,\golden^{\sigma(-1)}
    =3\,\golden^{-\lambda_{\log}}\,\golden^{-\log_{\golden}3}
    =3\,\golden^{-\lambda_{\log}}\cdot\tfrac13
    =\golden^{-\lambda_{\log}},
\]
so~(v) at $p=-1$ reads $\golden^{-\lambda_{\log}}=\tfrac12$.
\end{proof}

The factor $3=M_{2}$ (clause~(vii)) that divides out here is the same
multiplier the identity carries at every~$p$---the cancellation
$3\,\golden^{-\log_{\golden}3}=1$ is the mechanism of~(v), not a
coincidence of $p=-1$; what is proper to $p=-1$ is only that the binary
side reads $\tfrac12$.  The reciprocal turn therefore sits inside the
central identity~(v), the identity that carries the ternary factor
$3=M_{2}$.

\begin{Remark}[Inversion image of the rotor]\label{rem:inversion-image}
Clause~(i) inscribes $\golden$ as a root and its inverse,
$\golden=\zeta_{10}+\zeta_{10}^{-1}$.  Clause~(iv) carries the same
pattern---a value and its inverse---to the multiplicative line: the pair
$(\zeta_{10},\zeta_{10}^{-1})$ on the unit circle corresponds to
$(\golden^{\lambda_{\log}},\golden^{-\lambda_{\log}})=(2,\tfrac12)$
on~$\mathbb{R}_{>0}$.  The rotor is closed under inversion
($e^{i\theta}\leftrightarrow e^{-i\theta}$); the reciprocal turn
$\tfrac12$ is, on~$\mathbb{R}_{>0}$, the analogue of the conjugate root
$\zeta_{10}^{-1}$, completing the inversion symmetry of the
multiplicative companion so that it mirrors the rotor it accompanies.
\end{Remark}

These three inscriptions are internal: each holds by the machinery of
the present identity, with no operator invoked.  This same value
$\tfrac12$ is where the companion operator paper anchors a self-adjoint
operator in the Hilbert--P\'olya programme (\Cref{ssec:euler-substrate});
that operator and its relation to the zeros of~$\zeta$ are developed
there.

\section{Discussion}
\label{sec:discussion}

\subsection*{Identity and classification under equivalence}
The organising principle of modern and contemporary mathematics is
classification up to structural sameness: objects are known not in
themselves but through the isomorphisms that relate them.  This has a long
prehistory in the rhetorical structure of Vitruvius's \emph{De
Architectura}, whose opening dichotomy \emph{quod significatur et quod
significat}---the signified and the signifier, the \emph{res} and the
\emph{verba}---makes the work knowable through the proportional relations
its \emph{modulus} mediates: by number, proportion, and anatomy, in the
pre-Cartesian mode of interpretation Foucault
describes~\cite{santos-vitruvio,foucault-1966}.  The modulus is one unit
from which the proportions of the whole are read---identity generating
equivalence.

Modern mathematics inherits precisely this tension.  Two objects may be
\emph{equal}, literally the same, or merely \emph{isomorphic}, the same up
to a structure-preserving relabelling; classification proceeds almost always
by the second.  The present note is unusual in sitting at the meeting of the
two.  Its centre is a genuine identity---a single $\golden$, not a class of
isomorphic objects, satisfies twelve clauses at once---yet the clauses are
organised by equivalence: the residue classes of the Galois group
$\Ztwenty$, the split and inert types in $\mathbb{Q}(\sqrt5)$.  The same
$\golden$ is read through the equivalences that classify the structures it
inhabits.  The arc below moves between the two registers---identity where
the equation holds by construction, equivalence where the classification
does the work---answering the millennial question of
\Cref{ssec:classification} in the specific case of~$\golden$.

The note can be read as a single arc, from a tautology to a spectral
endpoint, each step naming where it is established.

\subsection*{The simultaneous setting}
The frame of the note is that a single $\golden$ is inscribed at once in
five canonical structures---the complex rotor of Euler's identity, the real
trigonometric line, the real multiplicative line, the arithmetic of Mersenne
modulo~$20$, and the Galois group of $\mathbb{Q}(\zeta_{20})/\mathbb{Q}$%
---the twelve simultaneous clauses of \Cref{thm:coincidence-intro}.  This is
the golden-ratio counterpart of Euler's identity: where $e^{i\pi}+1=0$ knots
analysis, algebra, and geometry at one point, here the same $\golden$ knots
those structures through the pentagon $\golden=2\cos(\pi/5)$ and its
cyclotomic inscription $\golden=\zeta_{10}+\zeta_{10}^{-1}$.  The central
identity below is $\golden$'s inscription on the multiplicative line; the
arc that follows traces its depth.

\subsection*{The tautological identity}
Read as an equality, the central identity $3\,\golden^{\sigma(p)}=2^{p}$ is
tautological: the golden index
$\sigma(p):=p\,\lambda_{\log}-\log_{\golden}3$ (\Cref{def:sigma}) is defined
precisely so that it holds, for every real~$p$, and its proof
(\Cref{thm:universal}) is the evaluation of that definition.  Nothing about
the truth of the equation is at stake; it holds by construction.  Yet the
slope it carries, $\lambda_{\log}=\log_{\golden}2$, is transcendental over
$\mathbb{Q}$ (Gelfond--Schneider, \Cref{rem:transcendence})---the exact
counterpart of the period $2\pi$ in Euler's identity, an algebraic relation
between algebraic numbers carried by a transcendental exponent.

\subsection*{Its non-tautological content}
What is not tautological is what the constituents of the equation turn out
to be.  The multiplier~$3$ is not a free constant: it equals
$M_{2}=|\{3,7,11\}|$, the number of residue classes the Mersenne primes
occupy modulo~$20$ (\Cref{thm:concentration-main},
\Cref{thm:three-main})---a fact about the arithmetic of $2^{p}-1$ that owes
nothing to the equation.  Here the identity first carries content beyond
its definition.

\subsection*{Its proof from outside}
That the three constituents---the factor~$3$, the base~$\golden$, the slope
$\lambda_{\log}$---are each pinned by a theorem \emph{external} to the
identity, and are mutually consistent with the golden coordinate uniquely
determined, is made precise in \Cref{thm:characterization}: the identity is
a consistency theorem for independently fixed determinations, certified
without circularity, not a discovered equality.  This is the sense in which
it is proved from outside itself.

\subsection*{Its expansion into the Mersenne primes}
Read forward, into the primes it produces, the identity reaches the 52
Mersenne exponents.  On the multiplicative line their binary powers
$2^{p}=\Mp+1$ fit a single integer sublattice of the golden base, every pair
standing in an exact golden ratio (\Cref{thm:ratio-main})---in the binary
guise, the shift of the single carried bit from one position to another
(\Cref{rem:binary})---with 52 distinct transcendental indices $\sigma(p_{i})$
(\Cref{cor:linear-independence}).

In the field $\mathbb{Q}(\sqrt5)$ the reach is finer, and it resolves each
Mersenne prime into presentations of one invariant.  The exponent residue
$p\bmod4$ fixes at once its class in $\Ztwenty$, its splitting type
(\Cref{cor:split-from-p}), and its Fibonacci sign $F_{\Mp}\equiv\pm1$
(\Cref{thm:fib-main}):
\[
  p\equiv1\mapsto(11,\ \mathrm{split},\ +1),\quad
  p\equiv3\mapsto(7,\ \mathrm{inert},\ -1),\quad
  p=2\mapsto(3,\ \mathrm{inert},\ -1),
\]
the triple recording the class modulo~$20$, the type in $\mathbb{Q}(\sqrt5)$,
and the sign of $F_{\Mp}$; and the inert primes are exactly those with
$\Mp\mid F_{\Mp+1}$ (\Cref{prop:rank-inert}).  Over the 52 known primes the
joint presentation takes exactly three values, one per admissible residue
(\Cref{cor:joint}); given $p\bmod4$, no further information separates the
class, the type, and the sign.  The development is \Cref{sec:fibonacci}.

\subsection*{The spectral meeting point}
Read backward, the same period $\lambda_{\log}$ carries $\golden$ down to
one half, $\golden^{-\lambda_{\log}}=\tfrac12$ (\Cref{sec:spectral}).  This
is the value at which the companion operator paper anchors its
Hilbert--P\'olya operator---the point at which the golden arithmetic
developed here meets the spectral programme.

\medskip
Throughout, the note works on the axis of representation and equivalence
rather than literal equality.  What a comparison of expressions registers as
distinct---$\Mp$ against $2^{p}$, an index defined against one proved---the
relating equivalences return to a single object in several guises.  This is
quantitative, not figurative: given the exponent residue $p\bmod4$, the
residue class, the splitting type, and the Fibonacci sign carry no
information beyond one another, so that over the 52 known Mersenne primes
their joint presentation takes exactly three values, one per admissible
residue.  The content lies not in any single presentation but in the exact
maps between them; where a literal reading sees distinct objects that fail to
coincide, the equivalence measures no difference at all.

\section{Formal verification in Lean 4}

The entire argument is formalised in a single self-contained file,
\texttt{golden\_mersenne.lean} (1599~lines, importing
only Mathlib, and including the Fibonacci splitting of~\Cref{sec:fibonacci}
and its consequences), built against Mathlib4 (Lean toolchain~\texttt{v4.28.0}, commit
\texttt{8f9d9cff6bd728b17a24e163c9402775d9e6a365}).  Gelfond--Schneider
is not present at this revision and enters as the explicit axiom
recorded below.  \Cref{tab:lean,tab:lean-b} list every Lean identifier
with its verification status, grouped by role.

\begin{table}[p]\centering\footnotesize
\setlength{\tabcolsep}{4pt}
\begin{tabular}{@{}p{1.0cm}p{3.6cm}p{5.0cm}p{1.5cm}@{}}
\toprule
Clause & Statement & Lean identifier & Status \\
\midrule

\multicolumn{4}{@{}l}{\emph{Clauses of \Cref{thm:coincidence-intro}}} \\
\midrule
(i)   & $\golden = \zeta_{10}+\zeta_{10}^{-1}$ &
  \lean{phi\_eq\_zeta10\_sum} & $0$ \texttt{sorry} \\
(ii)  & $\zeta_{10}^{10}=1$ &
  \lean{zeta10\_pow\_ten\_eq\_one} & $0$ \texttt{sorry} \\
(iii) & $\golden = 2\cos(\pi/5)$ &
  \lean{phi\_eq\_two\_cos\_pi\_fifth} & $0$ \texttt{sorry} \\
(iv)  & $\golden^{\lambda_{\log}}=2$ &
  \lean{mersenne\_bridge\_via\_lambda} & $0$ \texttt{sorry} \\
(v)   & $3\golden^{\sigma(p)}=2^{p}$ ($\forall p\in\mathbb{R}$) &
  \lean{golden\_tower\_bridge} & $0$ \texttt{sorry} \\
(vi)  & $\Mp\bmod 20\in\{3,7,11\}$ ($\forall p$ prime) &
  \lean{mersenne\_concentration\_general} & $0$ \texttt{sorry} \\
(vii) & $3 = M_{2} = |\{3,7,11\}|$ &
  \lean{three\_eq\_M2\_eq\_admissible\_card} & $0$ \texttt{sorry} \\
(ix)  & $|\Ztwenty|=8$ &
  \lean{ZtwentyStar\_card} & $0$ \texttt{sorry} \\
(viii), (x), (xi), (xii) & structural knottings articulated after~\Cref{thm:euler-connector} &
  \lean{phi\_euler\_mersenne\_connector} & $0$ \texttt{sorry} \\
\midrule
\multicolumn{4}{@{}l}{\emph{Auxiliary results (\Cref{sec:universal})}} \\
\midrule
--- & $\golden^{2}=\golden+1$ (defining relation) &
  \lean{phi\_sq} & $0$ \texttt{sorry} \\
--- & $\golden^{\log_{\golden} 3} = 3$ (\Cref{lem:ternary}) &
  \lean{phi\_pow\_log3} & $0$ \texttt{sorry} \\
--- & Uniqueness of $\sigma(p)$ (\Cref{thm:unique}) &
  \lean{sigma\_mersenne\_unique} & $0$ \texttt{sorry} \\
--- & $2^{q}/2^{p} = \golden^{(q-p)\lambda_{\log}}$ (\Cref{thm:ratio-main}) &
  \lean{mersenne\_ratio\_golden} & $0$ \texttt{sorry} \\
--- & 52-instance identity (\Cref{cor:51}) &
  \lean{mersenne\_phi\_correspondence\_52} & $0$ \texttt{sorry} \\
--- & Strict order of indices (\Cref{cor:51}) &
  \lean{sigma\_strictly\_increasing} & $0$ \texttt{sorry} \\
--- & Irrationality of $\lambda_{\log}$ &
  \lean{lambda\_log\_irrational} & $0$ \texttt{sorry} \\
--- & Transcendence of $\lambda_{\log}$ (\Cref{rem:transcendence}) &
  \lean{lambda\_log\_transcendental} & $0$ \texttt{sorry} \\
\midrule
\multicolumn{4}{@{}l}{\emph{Euler rotor (\Cref{eq:euler})}} \\
  \midrule
  --- & $e^{i\pi}=-1$ (rotor) &
    \lean{euler\_rotor} & $0$ \texttt{sorry} \\
\bottomrule
\end{tabular}
\caption{Lean~4 identifiers and verification status (part~I): clauses of
the simultaneous identities theorem, auxiliary results, and the Euler
rotor.}\label{tab:lean}
\end{table}

\begin{table}[p]\centering\footnotesize
\setlength{\tabcolsep}{4pt}
\begin{tabular}{@{}p{1.0cm}p{3.6cm}p{5.0cm}p{1.5cm}@{}}
\toprule
Clause & Statement & Lean identifier & Status \\
\midrule
\multicolumn{4}{@{}l}{\emph{Fibonacci splitting (\Cref{sec:fibonacci})}} \\
\midrule
--- & Binet identity in $\mathbb{Z}[\golden]$ &
  \lean{fib\_pow\_aux} & $0$ \texttt{sorry} \\
--- & Base cases $p\in\{2,3,5\}$ &
  \lean{mersenne\_fibonacci\_tabulated} & \texttt{native\allowbreak\_decide} \\
--- & Split/inert partition mod $5$ &
  \lean{chi5\_values} & $0$ \texttt{sorry} \\
--- & Lucas congruence $F_{q}\equiv(q/5)$ (\Cref{thm:fib-main}) &
  \lean{fibonacci\_splitting\_general} & $0$ \texttt{sorry} \\
--- & Split case (\Cref{lem:fib-split}) &
  \lean{fib\_split\_case} & $0$ \texttt{sorry} \\
--- & Inert case (\Cref{lem:fib-inert}) &
  \lean{fib\_inert\_case} & $0$ \texttt{sorry} \\
--- & Mersenne splitting type (\Cref{thm:residue-split}) &
  \lean{mersenne\_fibonacci\_type} & $0$ \texttt{sorry} \\
--- & Type from $p\bmod 4$ (\Cref{cor:split-from-p}) &
  \lean{mersenne\_mod5\_from\_p\_mod4} & $0$ \texttt{sorry} \\
--- & Rank characterises inert (\Cref{prop:rank-inert}) &
  \lean{mersenne\_rank\_iff\_inert} & $0$ \texttt{sorry} \\
--- & Joint presentation from $p\bmod 4$ (\Cref{cor:joint}) &
  \lean{mersenne\_presentation\allowbreak\_from\allowbreak\_p\_mod4} & $0$ \texttt{sorry} \\
\midrule
\multicolumn{4}{@{}l}{\emph{Identity characterization (\Cref{ssec:characterization})}} \\
\midrule
--- & Consistency \& uniqueness (\Cref{thm:characterization}) &
  \lean{golden\_identity\allowbreak\_characterization} & $0$ \texttt{sorry} \\
--- & Overdetermined consistency (\Cref{rem:characterization}) &
  \lean{golden\_identity\allowbreak\_is\allowbreak\_consistent\allowbreak\_overdetermined} & $0$ \texttt{sorry} \\
\midrule
\multicolumn{4}{@{}l}{\emph{External inputs (axioms)}} \\
\midrule
--- & Primality of $\Mp$, $p\leq 31$ &
  \lean{mersenne\_prime\_\{2..31\}} (8 thms) &
  \texttt{native\allowbreak\_decide} \\
--- & Primality of $\Mp$, $p\geq 61$ & 44 GIMPS \texttt{axiom}s &
  external \\
--- & Gelfond--Schneider (Hilbert's 7th) &
  \lean{transcendental\allowbreak\_cpow\allowbreak\_of\allowbreak\_isAlgebraic\allowbreak\_of\allowbreak\_irrational} & external \\
\bottomrule
\end{tabular}
\caption{Lean~4 identifiers and verification status (part~II): the
Fibonacci splitting, the identity characterization, and the external
inputs.}\label{tab:lean-b}
\end{table}

The external assumption set reduces to two items: (i)~the
Gelfond--Schneider theorem, declared as
\lean{transcendental\_cpow\_of\_isAlgebraic\_of\_irrational}, from which
the transcendence of $\lambda_{\log}$ over~$\mathbb{Q}$ is \emph{derived}
(\lean{lambda\_log\_transcendental}, \Cref{rem:transcendence}); and
(ii)~44 \texttt{axiom} declarations
$M_{61},\ldots,M_{82\,589\,933}\;\text{prime}$ (Lucas--Lehmer, GIMPS).
The golden relation $\golden^{2}=\golden+1$ and the irrationality of
$\lambda_{\log}$ are proved, not assumed. Gelfond--Schneider was
recently formalised in Lean~4~\cite{karatarakis-wiedijk} and is
being integrated into Mathlib; upon its landing the first axiom is
discharged. No \texttt{sorry} appears in any proof.

\section{Closing remark}

We have proved and formally verified a simultaneous identities
theorem for the golden ratio: a single $\golden$ satisfies
simultaneously twelve identities, inclusions, and identifications
across five canonical structures---the complex rotor of Euler's
identity, the real trigonometric line, the real-multiplicative line,
the arithmetic of Mersenne modulo~20, and the Galois group of the
cyclotomic extension $\mathbb{Q}(\zeta_{20})/\mathbb{Q}$.  The central
clause is the universal identity $3\,\golden^{\sigma(p)}=2^{p}$, a
mediation between the binary base of Mersenne arithmetic, the ternary
base $3=M_{2}$, and the golden base of pentagonal--cyclotomic
geometry through the bridge constant $\lambda_{\log}=\log_{\golden}2$.
The 52 currently known Mersenne prime exponents fit into this
structure as a distinguished integer sublattice of
$\golden^{\mathbb{Z}\lambda_{\log}-\log_{\golden} 3}$, and every pair
of Mersenne numbers admits an exact golden ratio.

\section*{Uses of generative artificial intelligence}

Generative AI tools were used by the authors as assistants in the
following parts of the preparation of this work: drafting and editing
the English prose; bibliographic search and cross-checking of
references; assistance in writing and debugging the Lean~4
formalisation; assistance in writing the numerical verification
scripts (mpmath/sympy); and organisation of the \LaTeX{} source.  All
mathematical results, definitions, and proofs are the authors' own,
and the entire manuscript has been extensively reviewed and corrected
by the human authors.  The correctness of every formal claim is
guaranteed by the Lean~4 kernel, which verifies each theorem against
its axioms independently of how the proof was produced; the
corresponding source compiles with $0$~\texttt{sorry}.  The authors
take full responsibility for the accuracy, originality, and integrity
of the manuscript.


\begin{thebibliography}{9}
\bibitem{boyer-1951}
C.~B.~Boyer,
\textit{The foremost textbook of modern times},
American Mathematical Monthly \textbf{58} (1951), no.~4, 223--226.

\bibitem{acevedo-2020}
Y.~A.~Acevedo Agudelo,
\textit{A complete classification of the Mersenne's primes and its
implications for computing},
Revista Politécnica \textbf{16} (2020), no.~32, 111--119.
DOI: 10.33571/rpolitec.v16n32a10.

\bibitem{acevedo-2021}
Y.~A.~Acevedo Agudelo,
\textit{Prime numbers: an alternative study using ova-angular
rotations},
preprint, arXiv:2104.04522, 2021.

\bibitem{connes-consani-site}
A.~Connes and C.~Consani,
\textit{The arithmetic site},
Comptes Rendus Mathématique \textbf{352} (2014), no.~12, 971--975;
and \textit{The scaling site},
Comptes Rendus Mathématique \textbf{354} (2016), no.~1, 1--6.

\bibitem{dirichlet}
P.~G.~L.~Dirichlet,
\textit{Beweis des Satzes, dass jede unbegrenzte arithmetische
Progression\ldots},
Abhandlungen der Königlich Preussischen Akademie der Wissenschaften
zu Berlin (1837), 45--81.

\bibitem{gelfond}
A.~O.~Gelfond,
\textit{Sur le septième problème de Hilbert},
Bulletin de l'Académie des Sciences de l'URSS \textbf{7} (1934), 623--634.

\bibitem{gimps}
G.~Woltman, S.~Kurowski et al.,
\textit{The Great Internet Mersenne Prime Search (GIMPS)},
\url{https://www.mersenne.org/primes/}.

\bibitem{godel-1931}
K.~Gödel,
\textit{Über formal unentscheidbare Sätze der Principia Mathematica
und verwandter Systeme~I},
Monatshefte für Mathematik und Physik \textbf{38} (1931), 173--198.

\bibitem{grisales}
J.~Grisales Herrera,
\textit{Golden Prime Symmetry},
B.Sc.\ thesis, Universidad del Tolima, Ibagué (Colombia),
November 2023.

\bibitem{kline-1972}
M.~Kline,
\textit{Mathematical Thought from Ancient to Modern Times},
Oxford University Press, New York, 1972.

\bibitem{lehmer}
D.~H.~Lehmer,
\textit{An extended theory of Lucas' functions},
Annals of Mathematics (2) \textbf{31} (1930), no.~3, 419--448.

\bibitem{karatarakis-wiedijk}
M.~Karatarakis and F.~Wiedijk,
\textit{A formalization of the Gelfond--Schneider theorem},
arXiv:2603.24823 (2026).

\bibitem{manin-1995}
Y.~I.~Manin,
\textit{Lectures on zeta functions and motives (according to Deninger
and Kurokawa)},
Astérisque \textbf{228} (1995), 121--163.

\bibitem{mathlib}
The Mathlib Community,
\textit{The Lean Mathematical Library},
Proceedings of CPP 2020, 367--381.

\bibitem{schneider}
T.~Schneider,
\textit{Transzendenzuntersuchungen periodischer Funktionen. I},
Journal für die reine und angewandte Mathematik \textbf{172} (1935), 65--69.

\bibitem{tarski-1948}
A.~Tarski,
\textit{A decision method for elementary algebra and geometry},
RAND Corporation Report R-109, Santa Monica, 1948; second edition,
University of California Press, Berkeley, 1951.

\bibitem{tits-1957}
J.~Tits,
\textit{Sur les analogues algébriques des groupes semi-simples
complexes},
in Colloque d'algèbre supérieure, Centre Belge de Recherches
Mathématiques, Établissement Ceuterick, Louvain, 1957, pp.~261--289.
\bibitem{santos-vitruvio}
G.~C. dos Santos,
\textit{Arquitetura vitruviana e ret\'orica antiga},
Archai \textbf{28} (2020), e02804.

\bibitem{foucault-1966}
M.~Foucault,
\textit{Les mots et les choses},
Gallimard, Paris, 1966.

\bibitem{weil-1940}
A.~Weil,
\textit{A 1940 letter to Simone Weil}, in
\textit{Collected Papers}, vol.~I, Springer, 1979, pp.~244--255.

\bibitem{weil-1948}
A.~Weil,
\textit{Sur les courbes alg\'ebriques et les vari\'et\'es qui s'en
d\'eduisent}, Hermann, Paris, 1948.

\end{thebibliography}

\end{document}
